% A Readymade beamer presentation template
% Version 1.1
% Relase date: May 2, 2010
% Released at http://www.stattler.com
% by Rifat Jahan

\documentclass{beamer}
%\usecolortheme[named=green]{structure}
\mode<presentation> {
\usetheme{Warsaw} % My favorite!
%\usetheme{Boadilla} % Pretty neat, soft color.
%\usetheme{default}
%\usetheme{Warsaw}
%\usetheme{Bergen} % This template has nagivation on the left
%\usetheme{Frankfurt} % Similar to the default with an extra region at the top.
%\usecolortheme{seahorse} % Simple and clean template
%\usetheme{Darmstadt} % not so good
% Uncomment the following line if you want page numbers and using Warsaw theme
% \setbeamertemplate{footline}[page number]
%\setbeamercovered{transparent}
\setbeamercovered{invisible}
% To remove the navigation symbols from the bottom of slides%
\setbeamertemplate{navigation symbols}{} 
}

\usepackage{tikz}
\usepackage{hyperref}
\usepackage{multicol}
\usepackage{animate}
\usepackage{caption} % to remove "Figure from caption
\usepackage{fancyvrb}
\usetikzlibrary{shapes,arrows}
\usepackage{graphicx}
\newcommand{\ac}{\`} % comando accento
\newcommand{\code}[1]{\texttt{#1}} 
\newcommand{\erpR}{\texttt{erpR }} 
\newcommand{\oq}{``} 
\newcommand{\cq}{''} 
\usepackage[backend=bibtex,style=verbose-trad2]{biblatex}

\newcommand{\btVFill}{\vskip0pt plus 1filll}



%\usepackage{Sweave} % sweave package

%% ------------- SECTION FOR MULTICOLUMNS TOC ------------%% 
% this code works only  if you also add \usepackage{multicol} above

\makeatletter
\def\insertsectionnavigation#1{%
  \hbox to #1{\vbox{{\usebeamerfont{section in head/foot}%
     \usebeamercolor[fg]{section in head/foot}%
     \def\slideentry##1##2##3##4##5##6{}%
     \def\sectionentry##1##2##3##4##5{%
       \ifnum##5=\c@part%
       \def\insertsectionhead{##2\hskip1em}%
       \def\insertsectionheadnumber{##1}%
       \def\insertpartheadnumber{##5}%
         \hyperlink{Navigation##3}{%
             \ifnum\c@section=##1%
               {\usebeamertemplate{section in head/foot}}%
             \else%
               {\usebeamertemplate{section in head/foot shaded}}%
             \fi%
         }\par
       \fi}%
       \parbox[c][0.25cm][c]{.5\paperwidth}{%
       \begin{multicols}{2}
       \dohead
       \end{multicols}}\space}
     }%
  \hfil}%
}

\def\insertsubsectionnavigation#1{%
  \hbox to #1{%
    \vbox{{%
      \usebeamerfont{subsection in head/foot}\usebeamercolor[fg]{subsection in head/foot}%
      \vskip0.5625ex%
      \beamer@currentsubsection=0%
      \def\sectionentry##1##2##3##4##5{}%
      \def\slideentry##1##2##3##4##5##6{\ifnum##6=\c@part\ifnum##1=\c@section%
        \ifnum##2>\beamer@currentsubsection%
        \beamer@currentsubsection=##2%
        \def\insertsubsectionhead{##5}%
        \def\insertsectionheadnumber{##1}%
        \def\insertsubsectionheadnumber{##2}%
        \def\insertpartheadnumber{##6}%
        \beamer@link(##4){%
              \ifnum\c@subsection=##2%
                {\usebeamertemplate{subsection in head/foot}}%
              \else%
                {\usebeamertemplate{subsection in head/foot shaded}}%
              \fi\hfill}\par
        \fi\fi\fi}%
       \hspace*{0.5em}\parbox[c][0cm][c]{\dimexpr.5\paperwidth-1em\relax}{%
       \begin{multicols}{2}
       \dohead\vskip0.5625ex\end{multicols}
       }\space
   }\hfil
}}}

\setbeamertemplate{headline}
{%
  \leavevmode\@tempdimb=2.4375ex%
  \ifnum\beamer@subsectionmax<\beamer@sectionmax%
    \multiply\@tempdimb by\beamer@sectionmax%
  \else%
    \multiply\@tempdimb by\beamer@subsectionmax%
  \fi%
  \ifdim\@tempdimb>0pt%
    \advance\@tempdimb by 1.125ex%
    \begin{beamercolorbox}[wd=.5\paperwidth,ht=0.5\@tempdimb,dp=2ex]{section in head/foot}%
      \vbox to0.5\@tempdimb{\vfill\insertsectionnavigation{.5\paperwidth}\vfill}%
    \end{beamercolorbox}%
    \begin{beamercolorbox}[wd=.5\paperwidth,ht=0.5\@tempdimb,dp=2ex]{subsection in head/foot}%
      \vbox to0.5\@tempdimb{\vfill\insertsubsectionnavigation{.5\paperwidth}\vfill}%
    \end{beamercolorbox}%
  \fi%
}
\makeatother

%% ------------- END OF SECTION FOR MULTICOLUMNS TOC ------------%% 



\captionsetup[figure]{labelformat=empty}% rede

%\usepackage{bm} 
% For typesetting bold math (not \mathbold)
%\logo{\includegraphics[height=0.6cm]{yourlogo.eps}}
%
 
\title[Rediscovering the Overlooked]{Rediscovering the overlooked in neuropsychology: \\ 
\vspace{0.2cm}
\small{how better methods can enhance clinical practice}}
%
\vfill
\author{\textbf{Giorgio Arcara}}
\vfill
\institute
{
IRCCS San Camillo Hospital, Venice, Italy\\
University of Padua, Padua, Italy
}
\date{February 18, 2025}

% \date \today \today will show current date. 
% Alternatively, you can specify a date.



\begin{document}

%
\begin{frame}
\titlepage
 \vspace*{0.5cm}
%\begin{figure}
\center{
%\includegraphics[scale=0.1]{Figures/XPRAGit_logo.png}\\
%\medskip
\href{https://creativecommons.org/licenses/by-nc-sa/4.0/}{\includegraphics[scale=0.15]{Figures/CC_BY-NC-SA.png}}\hspace{0.4em}{\tiny CC BY-NC-SA 4.0}\\
%\small{\url{https://creativecommons.org/licenses/by-nc-sa/3.0/}}
}
%end{figure}

\section{Introduction}

\end{frame}
%
%
% INTRODUCTION

\begin{frame}
\frametitle{A Personal Story}
\vfill
\begin{center}
\Large A personal story
\end{center}
%\pause
\begin{figure}
\fbox{\includegraphics[scale=0.25]{Figures/Tesi_Master.png}}
\caption*{\small 2006}
\end{figure}
\end{frame}


\begin{frame}
\frametitle{A personal Story}
\begin{columns}[c]
        \column{0.5\textwidth}
        \begin{figure}
        \includegraphics[width=0.6\textwidth]{Figures/Semeiotica.jpg}
        \caption*{\small 2016}
        \end{figure}
        \column{0.5\textwidth}
        %\pause
        \begin{figure}
        \includegraphics[width=0.6\textwidth]{Figures/Methodology.jpeg}
        \caption*{\small 2022}
        \end{figure}
\end{columns}
\end{frame}

\begin{frame}
\frametitle{The red thread}
\begin{center}
Can we improve the way we are doing clinical neuropsychology?
%\pause
\vfill
Do we need to create something new? Or do we need to discover something ``forgotten''? \\ %\pause
\vfill
\emph{(a step back to reflect on what we are doing)}
\end{center}
\end{frame}


\begin{frame}
\frametitle{A Personal Story}
\vfill
\begin{columns}
        \column{0.5\textwidth}
\begin{figure}
\fbox{\includegraphics[scale=0.25]{Figures/Tesi_Master.png}}
\caption*{\small 2006}
\end{figure}
        \column{0.5\textwidth}
\begin{itemize}
\item several missing pieces (no comprehensive manual on test development in neuropsychology).
\item attention on some aspects but not on others.
\end{itemize}
\end{columns}
\end{frame}

\begin{frame}
\frametitle{Back to the basics}
\begin{figure}
\includegraphics[scale=0.1]{Figures/Bookworm.png}
\end{figure}
Several books on Psychometrics later...
\vfill
It seems we have forgotten (and started to overlook) some basic properties of tests that need to be considered.
\end{frame}

\begin{frame}
\frametitle{Neuropsychological test in Italy}
\begin{center}
What is the current situation in Italy?
\end{center}
\end{frame}


\begin{frame}
\frametitle{Normative data}
It appears that Italian neuropsychological literature on tests is focused on \textbf{normative data}.
\vfill 
\begin{itemize}
\item Normative data are data collected on healthy participants to calculate thresholds to be used in clinical practice 
\item papers on normative data typically collect data on large samples ($N > 300$) and reports \emph{thresholds} and sometimes \emph{corrections}, to be applied (e.g. age and education)
\end{itemize}
\vfill
\begin{figure}
\includegraphics[scale=0.6]{Figures/Thresholds.png}
\end{figure}
\end{frame}



\begin{frame}
\frametitle{New tests in (Italian) Neuropsychology}
A PUBMED query
\begin{figure}
\includegraphics[scale=0.4]{Figures/NewTests.png}
\end{figure}
\texttt{\tiny{(((normative data[Title/Abstract]) OR ((norms[Title/Abstract]) OR ((normative values[Title/Abstract])))) AND (("2005/01/01"[Date - Publication] : "3000"[Date - Publication]))) 
AND (Test) AND (Italian [Title/Abstract]) NOT (Corrigendum OR Erratum OR Retract*[Title])}}
\end{frame}


\begin{frame}
\frametitle{The history of Italian methodology in neuropsychology}
\begin{columns}
        \column{0.5\textwidth}
\begin{figure}
\includegraphics[scale=0.3]{Figures/SpinnlerTognoni.png}
\caption*{\small 1987}
\end{figure}
        \column{0.5\textwidth}
        \begin{figure}
        \includegraphics[scale=0.2]{Figures/Bianchi.png}
        \caption*{\small 2008}
        \end{figure}
\end{columns}
\end{frame}

\begin{frame}
\frametitle{The history of Italian methodology in neuropsychology}
\begin{center}
Kit of the Italian Neuropsychologist
\begin{figure}
\includegraphics[scale=0.3]{Figures/Kit.png}
\caption*{\small 2011}
\end{figure}
\end{center}
\end{frame}



\begin{frame}
\emph{A speculation}
\vfill
I think that (in Italy) this is an \textbf{unwanted} consequence of the Spinnler \& Tognoni compendium (1987), which reported only normative data.
\vfill
It seems that this somehow was interperted as suggesting that only normative data are necessary to use a test {including newly developed ones}.
\end{frame}



\begin{frame}
\frametitle{A Systematic Review}
\begin{figure}
\includegraphics[scale=0.5]{Figures/Psychometrics.png}
\caption*{\href{https://link.springer.com/article/10.1007/s10072-021-05683-4}{\underline{Link}}}
\end{figure}
\end{frame}




\begin{frame}
\frametitle{Some questions}
\textbf{
\begin{itemize}
\item Are we certain that we are measuring (or assessing) what we intend to measure? %\pause
%\item Are we using correctly our tests when we monitor change over time?
\item Are age and education sufficient when assessing an individual? %\pause
\item Are we confident that we are using the best available methods to determine clinical thresholds?
\end{itemize}
}
\end{frame}





\section{Validity}

\begin{frame}
\begin{center}
{\Large{Are we certain that we are measuring (or assessing) what we intend to measure?}}
\vfill %\pause
\textbf{\Large Validity}
\end{center}
\end{frame}

\begin{frame}
Are we certain that our memory test is not heavily influenced by attention? Or executive functions?
\vfill
This has crucial implications on what we can conclude in clinical assessment. 
\end{frame}


\begin{frame}
\frametitle{Forgotten Validity}
A PUBMED query
\begin{figure}
\includegraphics[scale=0.4]{Figures/NewTests_validity.png}
\end{figure}
\tiny{\texttt{(((normative data[Title/Abstract]) OR ((norms[Title/Abstract]) OR ((normative values[Title/Abstract])))) AND (("2005/01/01"[Date - Publication] : "3000"[Date - Publication]))) 
AND (Test) AND (Italian [Title/Abstract])} \textbf{AND (Construct Validity)} \texttt{NOT (Corrigendum OR Erratum OR Retract*[Title])}}
\end{frame}




%%

\begin{frame}
\frametitle{Validity}
\textbf{Validity} is the quality of a test to measure what it intends to measure.
\vfill
%\pause
There are different ways to assess validity and different types of validity.
\begin{itemize}
\item construct validity
\item face validity
\item concurrent validity
\item predictive validity
\item ecological validity
\item \ldots
\end{itemize}
\end{frame}


\begin{frame}
\frametitle{Forgotten Validity}
In the review by Aiello et al. (2022) several tests used Concurrent Validity but not using \emph{Gold-Standards} with no thresholds for acceptable results.
\vfill
Almost \textbf{none} of the tests reported data on \textbf{Construct Validity*}.
\vfill 
\vfill
\vfill
*\footnotesize{sometimes it is not easy to understand whether some checks pertained construct validity}
\end{frame}

\begin{frame}
\frametitle{Forgotten Validity}
In the review by Aiello et al. (2022) several tests used Concurrent Validity but not using \emph{Gold-Standards} with no thresholds for acceptable results.
\vfill
Almost \textbf{none} of the tests reported data on \textbf{Construct Validity*}. \\
\vfill
\textbf{Are we really sure we are measuring what we aim to measure?}
\vfill
*\footnotesize{sometimes it is not easy to understand whether some checks pertained construct validity}
\end{frame}


\begin{frame}
\frametitle{Construct Validity}
How to assess construct validity? 
\vfill
In many ways, but mostly by inspecting if the test scores correlates (highly) with scores on test assessing the same construct.
\end{frame}



\begin{frame}
\frametitle{Rediscovered Validity}
\begin{figure}
\includegraphics[scale=0.4]{Figures/GEMS.png}
\end{figure}
\begin{center}
\href{https://onlinelibrary.wiley.com/doi/full/10.1002/brb3.2710}{\underline{Link}}
\end{center}
%
\end{frame}


\begin{frame}
\frametitle{GEMS Construct Validity}
\begin{itemize}
\item \underline{Correlation with MoCA = 0.72}  {\small(as in other articles, useful, but not sufficient!)}
\item \underline{Correlation with ENB-2 total score =  0.83} {\small(from a larger battery covering different aspects of cognition).}
\end{itemize}
%\pause
\vfill
Are we certain that ENB-2 total score is adequate? \\ (it is an arbitrary sum of scores)
\end{frame}


\begin{frame}
\frametitle{Tele-GEMS}
\begin{figure}
\includegraphics[scale=0.4]{Figures/TeleGEMS.png}
\end{figure}
\begin{center}
\href{https://link.springer.com/article/10.1007/s10072-023-06862-1}{\underline{Link}}
\end{center}
\end{frame}

\begin{frame}
\frametitle{A more sophisticated approach}
To develop a score on ``general cognitive ability'' on a group of patients in which validity was tested.
\vfill
\begin{itemize}
\item we take available scores on cognition \\ (SDMT, Rey's 15 words, and CET) %\pause
\item we perform a Principal Component Analysis (PCA) on the data. %\pause 
\item we extract the first Component explaining higher \% of variance (in this case it was 60\%).
\item we correlate Tele-GEMS with this bottom-up derived summary score.
\end{itemize}
%\pause
\underline{Correlation with Summary Score = 0.5}  \small{(significant but maybe not satisfactory).}
\end{frame}

\begin{frame}
\frametitle{Interim Conclusion (1/3)}
\vfill
We can (and should) investigate if we are actually measuring what we intend to measure.
\vfill
%\pause
In many cases we need to refine \emph{how} we can assess if the test is valid for its purpose (a job for methodologists).
\end{frame}


\section{Demographics}

\begin{frame}
\begin{center}
\item Are age and education sufficient when assessing an individual?
\vfill %\pause
\textbf{\Large Demographic corrections}
\end{center}
\end{frame}

%%%%

\begin{frame}
\frametitle{Test scores in clinical neuropsychology}
In a working memory test a patient scores 5.
\vfill
%\pause
In clinical neuropsychology, scores are almost always compared to \emph{thresholds} to facilitate interpretation.
\vfill
Typically, if the score is \emph{below} the threhsolds (i.e. \oq cut-offs \cq), it is intepreted as a probable impairment.
%\begin{block}
%\end{block}
\end{frame}
%

%%%
%\begin{frame}
%\frametitle{Demographic variables and normative data}
%Most methods for cognitive impairment are based on \emph{normative data}
%\vfill
%\begin{figure}
%\includegraphics[scale=0.18]{R_code_psicostat/Figures/normative_data.png}
%\end{figure}
%Threshold typically delimits 5\% of worse performances (at sample or population level)
%\end{frame}

\begin{frame}
\frametitle{Demographic variables and normative data}
Most methods for defining thresholds of cognitive impairment are based on \emph{normative data} that take into account age, education, and sex, using different methods.
\vfill
\begin{center}
\LARGE{ Why?}
\end{center}
\end{frame}

\begin{frame}
\frametitle{Demographic variables and normative data}
Age, education, and Sex are used in normative data because they are easy-to-use variables that are known to influence cognitive performance.
\vfill %\pause
\begin{center}
They tells us that, given those values in those variable, it is unlikely to observe that performance. \\
Hence we infer a probable cognitive impairment.
\end{center}
\vfill %\pause
In deriving our thresholds we may want to take into account these variables.
\end{frame}


\begin{frame}
\frametitle{Demographic variables and normative data}
 Demographic variables can be taken into acount by \textbf{dividing into bins}.
\vfill
\begin{figure}
\includegraphics[scale=0.2]{R_code_psicostat/Figures/normative_data_strat.png}
\end{figure}
\end{frame}
%
\begin{frame}
\frametitle{Demographic variables and normative data}
Or we can use \textbf{Regression Methods} that take into account the influence of age, education, and Sex.
\begin{figure}
\includegraphics[scale=0.5]{Figures/GEMS_demo.png}
\end{figure}
\end{frame}

\begin{frame}
\frametitle{Demographic variables and normative data}
Can we consider also other variables? Could they be relevant?
\vfill %\pause
This is an empirical question: if they helps us in estimating performance, than yes.
\end{frame}

\begin{frame}
\frametitle{MoCA and Cognitive Reserve}
\begin{figure}
\includegraphics[scale=0.4]{Figures/MOCA_CRIq.png}
\end{figure}
\begin{center}
\href{https://link.springer.com/article/10.1007/s12144-022-03062-6}{\underline{Link}}
\end{center}
\end{frame}

\begin{frame}
\frametitle{MoCA and Cognitive Reserve}
\textbf{Cognitive Reserve} is a consolidated concept in the field of aging and of cognitive performance.
\vfill
Several tools are available, as the \emph{Cognitive Reserve Index Questionnaire} (CRIq, Nucci et al., 2011)
\end{frame}

\begin{frame}
\frametitle{MoCA and Cognitive Reserve}
We tested several regression models, including or not CRIq, to check whether they were be useful to estimate performance.
\vfill %\pause
Models including CRIq as additional predictor yielded better results (measured with AIC).
\end{frame}

\begin{frame}
\frametitle{MoCA and Cognitive Reserve}
\begin{figure}
\includegraphics[scale=0.42]{Figures/MOCA_effects.png}
\end{figure}
\end{frame}



\begin{frame}
\frametitle{MoCA and Cognitive Reserve}
\begin{figure}
\includegraphics[scale=0.3]{Figures/MOCA_Shiny.png}
\end{figure}
\end{frame}

\begin{frame}
\frametitle{Ongoing results on Cognitive Reserve}
What about other neuropsychological tests?
%\pause
\vfill
\begin{figure}
\includegraphics[scale=0.4]{Figures/Biundo_ongoing.png}
\caption{(SINP, December 2024)}
\end{figure}
\end{frame}

\begin{frame}
\frametitle{Ongoing results on Cognitive Reserve}
In the battery \textbf{17/20 tests (85\%)}, performances were predicted more accurately when an index of Cognitive Reserve was included (Life Experience Questionnaire).
\end{frame}


\begin{frame}
\frametitle{Interim Conclusion (2/3)}
We should not forget why we are doing what we are doing (i.e., correcting for age and education).
\vfill %\pause
Maybe other variables can helps us in estimating cognitive performance, for more accurate clinical thresholds.
\end{frame}


%% 
\section{Methods}
\begin{frame}
\begin{center}
\item Are we confident that we are using the best available methods to determine clinical thresholds?
\vfill %\pause
\textbf{\Large Methods}
\end{center}
\end{frame}

\begin{frame}
\frametitle{Methodology in Clinical Neuropsychology}
Methodology in Clinical Neuropsychology should strive to improve the methods to yield more precise clinical data.
\vfill
%\pause
In Italy (see also my slides) we rely on a method with excellent rigour (that of Equivalent Scores), but are we sure this metod cannot be improved?
\end{frame}

%


\begin{frame}
\frametitle{Improving Equivalent Scores}
\begin{center}
\begin{figure}
\includegraphics[scale=0.4]{Figures/Improving.png}
\end{figure}
\href{https://link.springer.com/article/10.1007/s10072-024-07806-z}{\underline{Link}}
\end{center}
\end{frame}

\begin{frame}
\frametitle{Demographic variables and normative data}
Most methods for cognitive impairment are based on \emph{normative data} that take into account age, education, and sex, 
\begin{figure}
\includegraphics[scale=0.18]{R_code_psicostat/Figures/normative_data_cutoff.png}
\end{figure}
%\pause
Threshold typically delimits 5\% of worse performances (at sample or population level)
\end{frame}
%


\begin{frame}
\frametitle{Equivalent Scores}
\textbf{Equivalent Scores} is one of the most widespread method to calculate thresholds in clinical neuropsychology in Italy, originally developed by E. Capitani (Spinnler \& Tognoni, 1987).
%\pause
\vfill
the method consists of:
\begin{itemize}
\item a first part using regression to calculate adjusted scores that consider age, education, and sex.
\item a second part to derive thresholds for clinical inference (based on tolerance limits).
\end{itemize}
\end{frame}
%

%\begin{frame}
%\frametitle{Equivalent Scores}
%First calculate ES = 0.
%\begin{figure}
%\includegraphics[scale=0.6]{Figures/ES_1.png}
%\end{figure}
%\end{frame}
%
%
%\begin{frame}
%\frametitle{Equivalent Scores}
%Then calculate all other ES.
%\begin{figure}
%\includegraphics[scale=0.5]{Figures/ES_2.png}
%\end{figure}
%\end{frame}
%



%\begin{frame}
%\frametitle{Regression Method}
%\begin{figure}
%\includegraphics[scale=0.2]{R_code_psicostat/Figures/normative_data_example.png}
%\end{figure}
%This includes people of different Age (and Education, and Sex)
%\end{frame}
%


\begin{frame}
\frametitle{Regression Method}
\begin{figure}
\includegraphics[scale=0.25]{R_code_psicostat/Figures/age_effect.png}
\end{figure}
\end{frame}

\begin{frame}
\frametitle{Regression Method}
\begin{figure}
\includegraphics[scale=0.25]{R_code_psicostat/Figures/age_effect_lm.png}
\end{figure}
\end{frame}

\begin{frame}
\frametitle{Regression Method}
\begin{figure}
\includegraphics[scale=0.25]{R_code_psicostat/Figures/age_effect_resid1.png}
\end{figure}
\end{frame}

\begin{frame}
\frametitle{Regression Method}
\begin{figure}
\includegraphics[scale=0.25]{R_code_psicostat/Figures/age_effect_residall.png}
\end{figure}
\end{frame}

\begin{frame}
\frametitle{Regression Method}
\begin{figure}
\includegraphics[scale=0.25]{R_code_psicostat/Figures/hist_resid.png}
\end{figure}
* \footnotesize{Intercept is added to go back to original scale}
\end{frame}

\begin{frame}
\frametitle{Regression Method}
\begin{figure}
\includegraphics[scale=0.25]{R_code_psicostat/Figures/hist_resid_cutoff.png}
\end{figure}
* \footnotesize{Value is further adjusted so it goes back to original scale}
\end{frame}

\begin{frame}
\frametitle{Regression Modeling: two issues}
\begin{enumerate}
\item what predictors should be included in the model? (Age, Education, Sex)
\item what is the better transformation of predictors? (log, quadratic, inverse etc.)
\begin{figure}
\includegraphics[scale=0.25]{R_code_psicostat/Figures/transformations.png}
\end{figure}
\end{enumerate}
\end{frame}



\begin{frame}
\frametitle{Equivalent Scores - regression method \emph{(Capitani1987)}}
Capitani 1987 (Spinnler \& Tognoni 1987).
\begin{enumerate}
\item fit separate simple regression models for \emph{Age} and \emph{Education} trying different transformations from a list (e.g. log, inverse, quadratic, etc.)
\item fit a multiple regression with best transformations for Age and Education identified in 1. (best = higher $R^2$), and Sex.
\item drop terms with $p > 0.05$, Bonferroni corrected (i.e., $p < 0.017$)
\end{enumerate}
\end{frame}

\begin{frame}
\frametitle{An alternative method \emph{(Arcara2024)}}
Arcara 2024 (\emph{in prep.})
\begin{enumerate}
\item fit all possible multiple regressions including \emph{Age}, \emph{Education} and \emph{Sex} with all possible transformations for \emph{Age} and \emph{Education} (from a list). Calculate \textbf{AIC} for all models.
\item from each starting model in 1. , use a \emph{stepwise} backward selection (\code{step} function in \code{R}) to find the best model (lowest AIC) and select among all reduced models, select the model with the lowest AIC.
\item drop iteratevely terms with $p > 0.05$, Bonferroni corrected (i.e., $p < 0.05/number\ of\ terms$)
\end{enumerate}
\end{frame}

%\section{Method}


\begin{frame}
\frametitle{Comparison of methods}
Methods can be compared by exploting the fact that ES make explicit assumption about generative model of observed data, \underline{\textbf{simulating normative data}}
\vfill
\begin{itemize}
\item simulate normative data with known transformations and effects for Age, Education, and Sex
\item  Calculate how classifications made with ES calculated with methods overlap with \emph{Ground Truth} ES, (ES = 0, or ES $>$ 0)
\end{itemize}
\end{frame}

\begin{frame}
\frametitle{An example of simulation}
$19.5  -0.00031 \times Age^2   + 2.18 \times  log(Education)  + 0.258 \times Sex$
\begin{figure}
\includegraphics[scale=0.4]{Figures/sim_effects.png}
\end{figure}
\end{frame}

\begin{frame}
\frametitle{ES According to Ground Truth}
\center{Normative Data}
\begin{figure}
\includegraphics[scale=0.3]{Figures/dati_normativi_1.png}
\end{figure}
\end{frame}

\begin{frame}
\frametitle{ES According to Ground Truth}
\center{Normative Data}
\begin{figure}
\includegraphics[scale=0.3]{Figures/dati_normativi_2.png}
\end{figure}
%NOTE: nei dati normativi, ci si attende che  circa il 5\% delle performance (a livello di popolazione) saranno sotto cut-off.
\end{frame}

\begin{frame}
\frametitle{Performance Classification}
\begin{center}
\begin{tabular} {|c|c|c|}
 \hline
&  \textbf{GT ES = 0 } & \textbf{GT ES $>$ 0} \\ 
\hline
 \textbf{Method ES = 0 } & True Positive (TP) & False Positive (FP) \\ 
 \hline
\textbf{Method ES $>$ 0} & False Negative (FN) &  True Negative (TN)\\ 
 \hline
\end{tabular}
\end{center}
\vfill
%\pause
From these data we can derive
\begin{itemize}
\item  \textbf{Sensitivity} 
\item \textbf{Specificity}
\item \ldots
\end{itemize}
\end{frame}


%\begin{frame}
%\frametitle{Testing the methods}
%Methods can be compared by using the generative model assumed for ES and \underline{simulating normative data}.
%\vfill
%\begin{itemize}
%\item simulate data with known transformations of Age and Education and calculate \emph{Ground Truth} (GT) cut-offs.
%\end{itemize}
%\end{frame}

\begin{frame}
\frametitle{Parameters (1): Normative data Sample Size}
Normative Data Sample size\\
  $\begin{bmatrix} 100 & 300 & 500 & 700 \end{bmatrix}$
\end{frame}

\begin{frame}
\frametitle{Parameters (2): True Adjusted Score variability}
True Adjusted score Standard Deviation\\
  $\begin{bmatrix} 0.5 & 1 & 1.5  \end{bmatrix}$
  \begin{figure}
 \includegraphics[scale=0.25]{R_code_psicostat/Figures/True_adj_score_SD.png}
 \end{figure}
\end{frame}


%\begin{frame}
%\frametitle{Parameters (3): Relationship between Age and Education}
%\begin{itemize}
%\item independent distributions
%\item conditional distributions
%%\pause
%\begin{figure}
%\includegraphics[scale=0.4]{Figures/Age_edu_values.png}
%\end{figure}
%\end{itemize}
%\end{frame}

\begin{frame}
\frametitle{Simulation}
1000 simulation for each condition of the three parameters.\\
\vfill
\begin{itemize}
\item \emph{Normative data} (Train Set)
\item \emph{New participants} (i.e. Test Set of 1000 new participants for each Normative Data sample)
\end{itemize}
%\pause
\vfill
For each simulation performance is tested in classifying correctly below/above cut-off considering GT classifications (Sensitivity, Specificity, and others)
\vfill
NOTE: as data are simulated with healthy participants, about 0.05 (at population level) are expected to fall below cut-off.
\end{frame}

%\section{Results}

%\begin{frame}
%\frametitle{Mean Squared Error - Normative Data}
%\includegraphics[scale=0.18]{Figures/Performance_Av3_Cv2_sample.demo.cond_TRAIN_MSE.jpeg}
%\end{frame}

\begin{frame}
\frametitle{Sensitivity - Normative Data}
\includegraphics[scale=0.18]{Figures/Performance_Av3_Cv2_sample.demo.cond_TRAIN_SENSITIVITY.jpeg}
\end{frame}

\begin{frame}
\frametitle{Specificity - Normative Data}
\includegraphics[scale=0.18]{Figures/Performance_Av3_Cv2_sample.demo.cond_TRAIN_SPECIFICITY.jpeg}
\end{frame}

%\begin{frame}
%\frametitle{Precision - Normative Data}
%\includegraphics[scale=0.18]{Figures/Performance_Av3_Cv2_sample.demo.cond_TRAIN_PRECISION.jpeg}
%\end{frame}

%\begin{frame}
%\frametitle{F1 score - Normative Data}
%\includegraphics[scale=0.18]{Figures/Performance_Av3_Cv2_sample.demo.cond_TRAIN_F1.jpeg}
%\end{frame}

\begin{frame}
\frametitle{Sensitivity -  New Participants}
\includegraphics[scale=0.18]{Figures/Performance_Av3_Cv2_sample.demo.cond_TEST_SENSITIVITY.jpeg}
\end{frame}

\begin{frame}
\frametitle{Specificity - New Participants}
\includegraphics[scale=0.18]{Figures/Performance_Av3_Cv2_sample.demo.cond_TEST_SPECIFICITY.jpeg}
\end{frame}

%\begin{frame}
%\frametitle{Precision - New Participants}
%\includegraphics[scale=0.18]{Figures/Performance_Av3_Cv2_sample.demo.cond_TEST_PRECISION.jpeg}
%\end{frame}
%
%\begin{frame}
%\frametitle{F1 score - New Participants}
%\includegraphics[scale=0.18]{Figures/Performance_Av3_Cv2_sample.demo.cond_TEST_F1.jpeg}
%\end{frame}


\begin{frame}
\frametitle{Additional checks}
(not shown, but trust me)
\begin{itemize}
\item better performance in guessing correctly the Ground Truth transformation.
\item better performance with different conditional distributions of age and education.
\item better performance with different proportion of participants below/above cut-off.
\item the better method includes all three steps.
\end{itemize}
\end{frame}

%\section{Discussion}

\begin{frame}
\frametitle{Discussion}
\vfill
A simulation approach can help in develop better methods for clinical neuropsychology.
\vfill
The new proposed method outperforms the existing method across a wide range of parameters and conditions
\vfill
The Equivalent Score method (as many other methods in Clinical Neuropsychology) can be improved.
\end{frame} 


\begin{frame}
\frametitle{Future Directions}
Many aspects left:
\begin{itemize}
\item what about reliability of scores? (with Andrea Spoto)
\item what about heteroschedasticity and other realistic scenarios?
\item rising awareness into issues and limitations of cut-offs.
\end{itemize}
\end{frame} 

\begin{frame}
\frametitle{Resources}
\begin{itemize}
\item All code and scripts are available at \url{https://github.com/giorgioarcara/Reg-Method-Norm-Data-Sim}.\\
\item All supplementary materials available at \url{https://osf.io/yma69/}.
%\item Manuscript soon on Psyarxiv.
\item ShinyApp for easy application of the method {\small (with cautionary notes)}
\end{itemize}
\begin{figure}
\includegraphics[scale=0.1]{Figures/ShinyApp.png}
\end{figure}
\end{frame} 

\begin{frame}
\frametitle{Conclusions}
\begin{itemize}
\item We cannot take methods for granted, as they can always be improved to provide better tools for clinical practice. %\pause
\item Methodology allows us to take a step back and view things from fresh (or even traditional) but useful perspectives. %\pause
\item As researchers and clinicians, we should take the time to reflect: %\pause
\end{itemize}

\begin{center}
\textbf{Are we truly conducting the best possible clinical assessment?}
\end{center}

\end{frame}


\begin{frame}
\center
\Large{\textbf{Thanks!}}\\
\vfill
\normalsize{
giorgio.arcara@gmail.com\\
giorgio.arcara@hsancamillo.it\\
\url{giorgioarcara.github.io}
}

\end{frame} 



%\begin{frame}
%\frametitle{Multiple Linear Regressions}
%With \textbf{multiple} linear regression you can \emph{model} a response variable as a function of more than one predictor
%\begin{equation}
%Y =\beta_0 + \beta_1 X_1 + ... + \beta_n X_n + \epsilon
%\end{equation}
%\end{frame}




% End of slides
\end{document} 